\documentclass[12pt, a4paper]{article}

% ─── Paquetes ────────────────────────────────────────────────────────────────
\usepackage[utf8]{inputenc}
\usepackage[T1]{fontenc}
\usepackage[spanish]{babel}
\usepackage{amsmath, amssymb, amsthm}
\usepackage{geometry}
\usepackage{graphicx}
\usepackage{booktabs}
\usepackage{longtable}
\usepackage{array}
\usepackage{multirow}
\usepackage{caption}
\usepackage{subcaption}
\usepackage{float}
\usepackage{hyperref}
\usepackage[backend=biber, style=apa, sorting=nyt]{biblatex}
\usepackage{lmodern}
\usepackage{microtype}
\usepackage{setspace}
\usepackage{titlesec}
\usepackage{abstract}
\usepackage{fancyhdr}
\usepackage{xcolor}
\usepackage{tcolorbox}
\usepackage{siunitx}

\addbibresource{referencias.bib}

% ─── Geometría ───────────────────────────────────────────────────────────────
\geometry{
  a4paper,
  left=3cm, right=2.5cm,
  top=2.5cm, bottom=2.5cm
}

\onehalfspacing

% ─── Cabecera ────────────────────────────────────────────────────────────────
\pagestyle{fancy}
\fancyhf{}
\rhead{\small Prima ESG como factor Fama-French}
\lhead{\small Yan Ru Wu Jin}
\cfoot{\thepage}

% ─── Definiciones matemáticas ────────────────────────────────────────────────
\newtheorem{definition}{Definición}
\newtheorem{proposition}{Proposición}

% ─────────────────────────────────────────────────────────────────────────────
\begin{document}

% ─── Portada ─────────────────────────────────────────────────────────────────
\begin{titlepage}
  \centering
  \vspace*{1cm}
  {\large\textbf{UNIVERSIDAD COMPLUTENSE DE MADRID}\\[0.3em]
  Facultad de Ciencias Económicas y Empresariales}\\[2cm]

  \rule{\linewidth}{0.5mm}\\[0.5em]
  {\LARGE\bfseries La Prima de Mercado ESG como Factor Adicional\\[0.3em]
  en el Modelo de Tres Factores de Fama y French:\\[0.3em]
  Evidencia Empírica para el Mercado [\textit{ESPECIFICAR MERCADO}]}\\[0.3em]
  \rule{\linewidth}{0.5mm}\\[2cm]

  {\large Trabajo de Fin de Grado\\[0.5em]
  Grado en Economía / Administración y Dirección de Empresas}\\[2cm]

  \begin{tabular}{rl}
    \textbf{Autor:}    & Yan Ru Wu Jin \\[0.5em]
    \textbf{Director:} & [\textit{Nombre del director}] \\[0.5em]
    \textbf{Curso:}    & 2025--2026 \\
  \end{tabular}\\[3cm]

  \vfill
  {\large Madrid, junio de 2026}
\end{titlepage}

% ─── Resumen ─────────────────────────────────────────────────────────────────
\newpage
\begin{abstract}
\noindent
Este trabajo investiga si los factores ESG (Environmental, Social and Governance) generan
una prima de riesgo estadísticamente significativa que no es capturada por los modelos
de valoración de activos tradicionales. Para ello, se construye un factor ESG siguiendo
la metodología de Fama y French (1993), denominado \textit{HEL} (High ESG minus Low ESG),
y se contrasta su poder explicativo sobre los rendimientos en exceso de una muestra de
[\textit{N} empresas / índices] durante el período [\textit{fecha inicio}--\textit{fecha fin}].
Los resultados sugieren que [\textit{insertar resumen de hallazgos principales}].
El modelo de cuatro factores (FF3F + ESG) mejora el ajuste con respecto al modelo
de tres factores, con un $\bar{R}^2$ de [\textit{valor}] frente a [\textit{valor}].

\vspace{1em}
\noindent\textbf{Palabras clave:} ESG, prima de riesgo, modelo de Fama-French, inversión
sostenible, valoración de activos, factor HEL, CAPM.

\vspace{1em}
\noindent\textbf{Clasificación JEL:} G11, G12, G14, M14.
\end{abstract}

\newpage
\tableofcontents
\newpage

% ═════════════════════════════════════════════════════════════════════════════
\section{Introducción}
\label{sec:intro}
% ═════════════════════════════════════════════════════════════════════════════

En las últimas décadas, la inversión bajo criterios de sostenibilidad ha experimentado
un crecimiento exponencial. Según el \textit{Global Sustainable Investment Alliance}
(GSIA, 2022), los activos gestionados con algún criterio ESG superaron los 35 billones
de dólares en 2020, representando más del 35\,\% del total de activos bajo gestión
a nivel mundial. Este auge plantea una cuestión fundamental para la teoría financiera:
¿existe una \textit{prima de mercado ESG}? Es decir, ¿los inversores exigen una rentabilidad
adicional —o aceptan una inferior— por invertir en activos con mejores (o peores)
puntuaciones ESG?

La respuesta a esta pregunta tiene implicaciones tanto teóricas como prácticas.
Desde el punto de vista teórico, cuestiona si los modelos estándar de valoración
de activos —como el CAPM de Sharpe (1964) y Lintner (1965) o el modelo de tres
factores de Fama y French (1993)— son suficientes para explicar la sección transversal
de rendimientos, o si existe una fuente de riesgo sistemático vinculada a la
sostenibilidad que estos modelos omiten. Desde el punto de vista práctico, afecta
a las decisiones de asignación de activos, la gestión de carteras responsables y
el diseño de productos financieros ESG.

La literatura académica no ha llegado a un consenso. Mientras que algunos trabajos
encuentran que las empresas con mejores puntuaciones ESG obtienen rendimientos superiores
ajustados por riesgo (Edmans, 2011; Eccles, Ioannou y Serafeim, 2014), otros
documentan que los inversores ``verdes'' aceptan rendimientos inferiores como
compensación no financiera (Pástor, Stambaugh y Taylor, 2021, 2022), y un tercer
grupo no encuentra relación estadísticamente significativa (Hong y Kacperczyk, 2009).

Este trabajo contribuye a ese debate mediante la construcción de un factor ESG
adicional —que denominamos \textit{HEL} (High ESG minus Low ESG)— siguiendo
la metodología de construcción de carteras de Fama y French (1993), e incorporándolo
al modelo de tres factores para obtener un modelo de cuatro factores.
[\textit{Describir brevemente los datos y el mercado analizado.}]

El resto del trabajo se organiza de la siguiente manera: la Sección~\ref{sec:marco}
presenta el marco teórico; la Sección~\ref{sec:revision} revisa la literatura
empírica relevante; la Sección~\ref{sec:metodologia} describe la metodología;
la Sección~\ref{sec:datos} detalla los datos; la Sección~\ref{sec:resultados}
presenta los resultados; y la Sección~\ref{sec:conclusiones} concluye.

% ═════════════════════════════════════════════════════════════════════════════
\section{Marco Teórico}
\label{sec:marco}
% ═════════════════════════════════════════════════════════════════════════════

\subsection{El modelo CAPM}
\label{subsec:capm}

El modelo de valoración de activos de capital (\textit{Capital Asset Pricing Model},
CAPM) fue desarrollado de forma independiente por Sharpe (1964), Lintner (1965)
y Mossin (1966) a partir de la teoría de selección de carteras de Markowitz (1952).
Bajo los supuestos de mercados perfectos, expectativas homogéneas e inversores
aversos al riesgo con función de utilidad de varianza media, el CAPM establece que
la rentabilidad esperada de cualquier activo $i$ satisface:

\begin{equation}
  \label{eq:capm}
  \mathbb{E}[R_i] - R_f = \beta_i \, \bigl(\mathbb{E}[R_m] - R_f\bigr),
\end{equation}

donde $R_i$ es la rentabilidad del activo $i$, $R_f$ es la tasa libre de riesgo,
$R_m$ es la rentabilidad de la cartera de mercado y $\beta_i$ es la covarianza de $R_i$
con $R_m$ normalizada por la varianza de $R_m$:

\begin{equation}
  \label{eq:beta}
  \beta_i = \frac{\mathrm{Cov}(R_i, R_m)}{\mathrm{Var}(R_m)}.
\end{equation}

La magnitud $\mathbb{E}[R_m] - R_f$ se denomina \textit{prima de mercado} y constituye
el precio del riesgo sistemático en el CAPM. En su versión empírica, el CAPM se
estima mediante la regresión de series temporales propuesta por Black, Jensen y
Scholes (1972):

\begin{equation}
  \label{eq:capm_ts}
  R_{i,t} - R_{f,t} = \alpha_i + \beta_i \, (R_{m,t} - R_{f,t}) + \varepsilon_{i,t},
\end{equation}

donde $\alpha_i$ (``alfa de Jensen'') mide la rentabilidad anormal no explicada
por el factor de mercado. Si el CAPM es válido, $\alpha_i = 0$ para todos los activos.

Sin embargo, una extensa literatura empírica ha documentado la existencia de
\textit{anomalías} —patrones sistemáticos de rentabilidad no explicados por el CAPM—
que constituyen evidencia en su contra. Las más robustas incluyen el \textit{efecto tamaño}
(Banz, 1981), el \textit{efecto valor} (Rosenberg, Reid y Lanstein, 1985;
Fama y French, 1992) y el efecto \textit{momentum} (Jegadeesh y Titman, 1993).

\subsection{El modelo de tres factores de Fama y French}
\label{subsec:ff3f}

Para capturar las anomalías de tamaño y valor, Fama y French (1993) proponen ampliar
el CAPM con dos factores adicionales. Su modelo de tres factores (FF3F) es:

\begin{equation}
  \label{eq:ff3f}
  R_{i,t} - R_{f,t} = \alpha_i + \beta_i \, \mathrm{MKT}_t
    + s_i \, \mathrm{SMB}_t + h_i \, \mathrm{HML}_t + \varepsilon_{i,t},
\end{equation}

donde:

\begin{itemize}
  \item $\mathrm{MKT}_t = R_{m,t} - R_{f,t}$ es el factor de mercado (prima de riesgo
        de mercado).

  \item $\mathrm{SMB}_t$ (\textit{Small Minus Big}) captura la prima de tamaño: es
        la diferencia en rentabilidad entre carteras de empresas pequeñas y grandes.

  \item $\mathrm{HML}_t$ (\textit{High Minus Low}) captura la prima de valor: es la
        diferencia en rentabilidad entre carteras de empresas con alto y bajo ratio
        valor contable/valor de mercado (\textit{book-to-market}).
\end{itemize}

La construcción de SMB y HML sigue una metodología basada en intersecciones de carteras.
En junio de cada año $t$, Fama y French ordenan las empresas por capitalización bursátil
en dos grupos: \textit{Small} (S) y \textit{Big} (B), usando la mediana como punto de
corte. Simultáneamente, las ordenan por su ratio \textit{book-to-market} en tres grupos:
\textit{Low} (L, percentil 0--30), \textit{Medium} (M, percentil 30--70) y
\textit{High} (H, percentil 70--100). La intersección de estas clasificaciones genera
seis carteras: SL, SM, SH, BL, BM, BH.

\begin{equation}
  \label{eq:smb}
  \mathrm{SMB}_t = \frac{1}{3}(R_{SL,t} + R_{SM,t} + R_{SH,t})
    - \frac{1}{3}(R_{BL,t} + R_{BM,t} + R_{BH,t}),
\end{equation}

\begin{equation}
  \label{eq:hml}
  \mathrm{HML}_t = \frac{1}{2}(R_{SH,t} + R_{BH,t})
    - \frac{1}{2}(R_{SL,t} + R_{BL,t}).
\end{equation}

Fama y French (1993) demuestran que este modelo captura buena parte de la variación
en los rendimientos de las carteras de acciones estadounidenses, con $R^2$ cercanos
al 90--95\,\% en carteras bien diversificadas. Los coeficientes $s_i > 0$ y $h_i > 0$
indican que el activo $i$ se comporta más como empresas pequeñas y de valor,
respectivamente.

\subsection{Extensiones del modelo: el factor de cuatro y cinco factores}
\label{subsec:extensiones}

La literatura ha propuesto numerosas extensiones al FF3F. Las más influyentes son:

\begin{enumerate}
  \item \textbf{Modelo de cuatro factores de Carhart (1997)}: añade un factor
        \textit{momentum} ($\mathrm{WML}_t$, \textit{Winners Minus Losers}) basado
        en la diferencia de rentabilidad entre acciones ganadoras y perdedoras en
        los últimos 12 meses.

  \item \textbf{Modelo de cinco factores de Fama y French (2015)}: incorpora
        dos factores adicionales —\textit{RMW} (\textit{Robust Minus Weak},
        prima de rentabilidad) y \textit{CMA} (\textit{Conservative Minus Aggressive},
        prima de inversión)— basados en la teoría de valoración por dividendos de
        Miller y Modigliani.
\end{enumerate}

La lógica subyacente a estas extensiones es siempre la misma: si existe un patrón
sistemático de rendimientos que el modelo vigente no explica, puede deberse a un
factor de riesgo omitido. Este es precisamente el planteamiento que motiva la
incorporación del factor ESG.

\subsection{Fundamentos teóricos de la prima ESG}
\label{subsec:teoria_esg}

Existen al menos tres marcos teóricos que ofrecen predicciones sobre la relación
entre criterios ESG y rendimientos financieros:

\subsubsection{Teoría de la compensación por riesgo}

Si el bajo desempeño ESG (por ejemplo, elevadas emisiones de carbono, escándalos
de gobernanza o prácticas laborales cuestionables) constituye un riesgo sistemático
—es decir, no diversificable—, los inversores exigirán una mayor rentabilidad
esperada para compensarlo. En este caso, las empresas con menores puntuaciones
ESG deberían ofrecer mayores rendimientos esperados (\textit{prima ESG positiva}
para la cartera Long Low ESG - Short High ESG). Formalmente, el factor HEL tendría
un precio negativo: las empresas muy ESG generan menos rentabilidad esperada porque
son activos de cobertura frente al riesgo de sostenibilidad.

\subsubsection{Teoría de la utilidad no pecuniaria (\textit{taste-based preferences})}

Pástor, Stambaugh y Taylor (2021, 2022) desarrollan un modelo de equilibrio en el
que algunos inversores obtienen utilidad directa (\textit{no pecuniaria}) de invertir
en activos verdes. En equilibrio, esto eleva el precio de los activos ESG, reduciendo
su rendimiento esperado futuro. Sin embargo, un aumento inesperado en las preferencias
de los inversores por activos ESG puede generar rendimientos \textit{realizados}
superiores en el corto plazo, aunque los rendimientos \textit{esperados} sean inferiores.
Este marco reconcilia la aparente contradicción entre los hallazgos empíricos.

\subsubsection{Teoría de la información y los mercados eficientes}

Bajo la hipótesis de mercados eficientes (Fama, 1970), si la información ESG está
completamente reflejada en precios, no debería existir una prima ESG sistemática
explotable (neta de costes de transacción). No obstante, si existen fricciones
de mercado —como la inversión de muchos fondos institucionales en universos ESG
restringidos— puede surgir una prima de liquidez o de exclusión.

% ═════════════════════════════════════════════════════════════════════════════
\section{Revisión de la Literatura}
\label{sec:revision}
% ═════════════════════════════════════════════════════════════════════════════

\subsection{ESG y rendimientos financieros: evidencia general}
\label{subsec:lit_general}

La relación entre ESG y rendimiento financiero ha sido estudiada extensamente.
Friede, Busch y Bassen (2015) realizaron un meta-análisis de más de 2.200 estudios
publicados entre 1970 y 2014, concluyendo que aproximadamente el 90\,\% de ellos
encuentran una relación no negativa entre ESG y rendimiento corporativo. No obstante,
el metaanálisis también revela una importante heterogeneidad en los métodos y en las
definiciones de ESG empleadas.

En la literatura de asset pricing, Hong y Kacperczyk (2009) documentan que las
denominadas ``\textit{sin stocks}'' (tabaco, alcohol, armas, juego) obtienen
rendimientos ajustados por riesgo superiores, consistente con una \textit{prima
de exclusión} exigida por inversores con restricciones ESG. Esto implica una
prima positiva para las carteras de bajo ESG.

Por el contrario, Edmans (2011) encuentra que las empresas en la lista ``100 Best
Companies to Work For'' de Fortune generan rendimientos anormales positivos durante
1984--2009, incluso controlando por factores de riesgo estándar. Eccles, Ioannou
y Serafeim (2014) hallan resultados similares utilizando un índice de sostenibilidad
compuesto.

\subsection{ESG como factor de riesgo sistemático}
\label{subsec:lit_factor}

La cuestión de si los criterios ESG constituyen un factor de riesgo sistemático
ha motivado trabajos que adoptan directamente la metodología de Fama y French.

Pedersen, Fitzgibbons y Pomorski (2021) proponen un modelo de equilibrio donde
los inversores difieren en su acceso a información ESG. En su framework, las
puntuaciones ESG afectan tanto al riesgo percibido como a los rendimientos esperados,
y demuestran que una \textit{frontera eficiente ESG} domina la frontera media-varianza
estándar. Su validación empírica muestra que los activos con altas puntuaciones
ESG tienen betas más bajas y menores costes de capital.

Pástor, Stambaugh y Taylor (2021) desarrollan el marco teórico más completo
hasta la fecha. En su modelo, los inversores heterogéneos difieren en sus
``gustos'' por activos verdes. Las acciones verdes tienen menor rendimiento esperado
en equilibrio (su precio incorpora la utilidad no pecuniaria). Sin embargo, cuando
el gusto agregado por la sostenibilidad crece de forma inesperada, las acciones
verdes obtienen rendimientos realizados superiores. Esto genera una distinción
crucial entre rendimientos \textit{esperados} (negativamente relacionados con ESG)
y rendimientos \textit{realizados} (potencialmente positivos en periodos de auge ESG).

En su trabajo posterior (Pástor, Stambaugh y Taylor, 2022), estiman este modelo
con datos de fondos de inversión sostenible y encuentran evidencia empírica consistente:
los fondos ESG obtuvieron rendimientos inferiores en el largo plazo, pero rendimientos
superiores durante el periodo de creciente concienciación ESG (2016--2020).

\subsection{Construcción de factores ESG tipo Fama-French}
\label{subsec:lit_metodologia}

Varios trabajos han construido factores ESG siguiendo la metodología de carteras
ordenadas de Fama y French.

Auer y Schuhmacher (2016) construyen carteras long-short basadas en puntuaciones
ESG para el mercado europeo, encontrando que las diferencias de rendimiento no
son estadísticamente significativas para la mayoría de industrias y subsegmentos.

Nagy, Kassam y Lee (2016) utilizan datos de MSCI ESG Ratings para construir
un factor ESG para el mercado global, encontrando que la integración ESG mejora
el perfil riesgo-rendimiento de las carteras, aunque el efecto varía
considerablemente entre regiones.

Lioui y Tarelli (2022) demuestran que el factor ESG carga sobre los factores
de FF3F (especialmente sobre MKT y HML), lo que sugiere que parte de la prima ESG
puede ser redundante. Sin embargo, su carga sobre el factor $\alpha$ (alfas anormales)
es significativa en algunos submercados.

\subsection{Debate abierto: ¿prima positiva o negativa?}
\label{subsec:lit_debate}

El debate sobre el signo de la prima ESG no está cerrado. La evidencia puede
sintetizarse de la siguiente manera:

\begin{itemize}
  \item \textbf{Prima ESG positiva} (High ESG $>$ Low ESG): consistente con la
        hipótesis de que el buen desempeño ESG reduce el riesgo idiosincrático,
        aumenta la calidad crediticia o genera externalidades positivas capturadas
        en los beneficios a largo plazo.

  \item \textbf{Prima ESG negativa} (Low ESG $>$ High ESG): consistente con la
        hipótesis de que los inversores ESG aceptan menor rentabilidad financiera
        como compensación por la utilidad no pecuniaria, o con la existencia de
        una prima de exclusión para activos ``pecaminosos''.

  \item \textbf{Prima ESG nula}: consistente con mercados eficientes donde la
        información ESG ya está incorporada en precios, o con la heterogeneidad
        de puntuaciones entre proveedores que dificulta la identificación empírica.
\end{itemize}

[\textit{Posicionar los resultados propios en este debate.}]

% ═════════════════════════════════════════════════════════════════════════════
\section{Metodología}
\label{sec:metodologia}
% ═════════════════════════════════════════════════════════════════════════════

\subsection{El modelo de cuatro factores propuesto}
\label{subsec:modelo}

Siguiendo la lógica de Fama y French (1993), se propone un modelo de cuatro factores
que incorpora un factor ESG:

\begin{equation}
  \label{eq:ff4f}
  R_{i,t} - R_{f,t} = \alpha_i + \beta_i \, \mathrm{MKT}_t
    + s_i \, \mathrm{SMB}_t + h_i \, \mathrm{HML}_t
    + e_i \, \mathrm{HEL}_t + \varepsilon_{i,t},
\end{equation}

donde $\mathrm{HEL}_t$ (\textit{High ESG minus Low ESG}) es el nuevo factor de
sostenibilidad, y $e_i$ es su carga factorial para el activo $i$.

Un coeficiente $e_i > 0$ indica que el activo se comporta como activos de alta
puntuación ESG; $e_i < 0$ indica comportamiento similar a activos de baja puntuación ESG.

El parámetro de interés central es el \textit{precio del factor HEL}: si el factor
HEL tiene un rendimiento esperado significativamente distinto de cero, implica que
los inversores exigen (o aceptan) una prima por exposición al riesgo ESG.

\subsection{Construcción del factor HEL}
\label{subsec:hel}

La construcción del factor HEL sigue la metodología estándar de Fama y French:

\paragraph{Clasificación ESG.}
En cada período de rebalanceo, se ordenan las empresas de la muestra según su
puntuación ESG [\textit{especificar fuente: MSCI, Sustainalytics, Bloomberg, etc.}].
Se forman tres grupos: \textit{High} (H, percentil 0--[\textit{X}]\%), \textit{Medium} (M)
y \textit{Low} (L, percentil [\textit{Y}]--100\%). [\textit{Describir puntos de corte usados.}]

\paragraph{Clasificación por tamaño.}
Simultáneamente, las empresas se dividen en dos grupos por capitalización bursátil:
\textit{Small} (S) y \textit{Big} (B), usando la mediana como punto de corte.

\paragraph{Carteras de intersección.}
La intersección de las clasificaciones ESG y tamaño genera seis carteras:
SH, SM, SL, BH, BM, BL. Cada cartera se pondera por capitalización bursátil (\textit{value-weighted}).

\paragraph{Cálculo del factor HEL.}
\begin{equation}
  \label{eq:hel}
  \mathrm{HEL}_t = \frac{1}{2}(R_{SH,t} + R_{BH,t})
    - \frac{1}{2}(R_{SL,t} + R_{BL,t}).
\end{equation}

Este factor mide la rentabilidad diferencial de las empresas con alta puntuación ESG
frente a las de baja puntuación, controlando por el efecto tamaño.

\subsection{Estimación y contraste}
\label{subsec:estimacion}

Los modelos se estiman por \textit{Mínimos Cuadrados Ordinarios} (MCO) en series
temporales para cada activo (o cartera) $i$:

\begin{enumerate}
  \item \textbf{Modelo CAPM} [ecuación~(\ref{eq:capm_ts})]: regresión del exceso de
        rentabilidad sobre el factor de mercado.

  \item \textbf{Modelo FF3F} [ecuación~(\ref{eq:ff3f})]: regresión sobre MKT, SMB y HML.

  \item \textbf{Modelo FF4F-ESG} [ecuación~(\ref{eq:ff4f})]: regresión sobre MKT, SMB,
        HML y HEL.
\end{enumerate}

Para evaluar la significatividad del factor ESG se utiliza el test $t$ del coeficiente
$e_i$ con errores estándar robustos (Newey-West con [\textit{número} lags] para corregir
autocorrelación y heterocedasticidad). Para comparar los modelos anidados se emplea
el test $F$ de significatividad conjunta de los factores añadidos.

La bondad de ajuste se mide mediante el $\bar{R}^2$ ajustado. La capacidad del
modelo para explicar la sección transversal de rendimientos se evalúa mediante:

\begin{itemize}
  \item Los alfas $\hat\alpha_i$: un modelo bien especificado debería producir alfas
        estadísticamente no significativos.
  \item El estadístico $GRS$ de Gibbons, Ross y Shanken (1989), que contrasta
        conjuntamente la hipótesis $H_0: \alpha_1 = \cdots = \alpha_N = 0$.
\end{itemize}

% ═════════════════════════════════════════════════════════════════════════════
\section{Datos}
\label{sec:datos}
% ═════════════════════════════════════════════════════════════════════════════

\subsection{Muestra de empresas y período de análisis}
\label{subsec:muestra}

[\textit{Describir: número de empresas, mercado(s) analizado(s), período temporal,
criterios de inclusión/exclusión (mínimo de datos disponibles, exclusión de sectores
financieros si procede, etc.).}]

La muestra está compuesta por [\textit{N}] empresas pertenecientes a
[\textit{índice o universo, p.ej. IBEX 35, STOXX 600, S\&P 500}]
durante el período [\textit{fecha inicio}]--[\textit{fecha fin}],
con datos [\textit{frecuencia: mensual / semanal}].
Se excluyen [\textit{criterios de exclusión}].

\subsection{Datos de rentabilidades}
\label{subsec:datos_rentabilidades}

Los datos de precios de cierre ajustados provienen de
[\textit{fuente: Yahoo Finance / Bloomberg / Refinitiv}].
La rentabilidad simple del activo $i$ en el período $t$ se calcula como:
\[
  R_{i,t} = \frac{P_{i,t} - P_{i,t-1} + D_{i,t}}{P_{i,t-1}},
\]
donde $P_{i,t}$ es el precio ajustado y $D_{i,t}$ son los dividendos.
Como tasa libre de riesgo se utiliza [\textit{especificar: Letras del Tesoro
a 3 meses / Euribor / tasa de la Fed Funds}].

Los factores MKT, SMB y HML provienen de
[\textit{fuente: Kenneth French Data Library / elaboración propia}].

\subsection{Puntuaciones ESG}
\label{subsec:datos_esg}

Las puntuaciones ESG se obtienen de [\textit{fuente: MSCI ESG Ratings /
Sustainalytics / Bloomberg ESG Disclosure Score / Refinitiv ESG}].
[\textit{Describir brevemente la metodología de la fuente y las características
de las puntuaciones: rango, frecuencia de actualización, cobertura.}]

Para evitar el sesgo de \textit{look-ahead}, las puntuaciones ESG de cada año $t$
se asignan a partir de julio de $t+1$ (cuando están públicamente disponibles).

\subsection{Estadísticos descriptivos}
\label{subsec:estadisticos}

El Cuadro~\ref{tab:descriptivos} presenta los estadísticos descriptivos de
los factores y de los excesos de rentabilidad de las carteras.

\begin{table}[H]
  \centering
  \caption{Estadísticos descriptivos de los factores ([\textit{periodo}])}
  \label{tab:descriptivos}
  \begin{tabular}{lrrrrrrr}
    \toprule
    & Media & Mediana & Std. Dev. & Mín. & Máx. & Asimetría & Curtosis \\
    \midrule
    MKT     & \multicolumn{7}{c}{[\textit{rellenar con datos propios}]} \\
    SMB     & & & & & & & \\
    HML     & & & & & & & \\
    HEL     & & & & & & & \\
    $R_f$   & & & & & & & \\
    \bottomrule
  \end{tabular}
  \begin{minipage}{\linewidth}
    \small\vspace{0.3em}
    \textit{Nota:} Todos los valores en \%. [\textit{Frecuencia temporal.}]
    El factor MKT es la prima de riesgo de mercado $(R_m - R_f)$.
    HEL es el factor ESG construido siguiendo la metodología descrita
    en la Sección~\ref{subsec:hel}.
  \end{minipage}
\end{table}

% ═════════════════════════════════════════════════════════════════════════════
\section{Resultados}
\label{sec:resultados}
% ═════════════════════════════════════════════════════════════════════════════

\subsection{Evolución temporal del factor HEL}
\label{subsec:hel_evolucion}

[\textit{Describir la evolución temporal del factor HEL. Incluir gráfico.}]

La Figura~\ref{fig:hel_evolucion} muestra la evolución del factor HEL acumulado
durante el período de estudio.

\begin{figure}[H]
  \centering
  % \includegraphics[width=0.85\textwidth]{../figures/hel_evolucion.png}
  \fbox{\parbox{0.85\textwidth}{\centering\vspace{3em}[Figura: Evolución del factor HEL acumulado]\vspace{3em}}}
  \caption{Rentabilidad acumulada del factor HEL (\textit{High ESG minus Low ESG}).}
  \label{fig:hel_evolucion}
\end{figure}

[\textit{Análisis descriptivo: ¿El factor HEL es positivo o negativo en el período?
¿Cuál es su prima media anualizada? ¿Hay subperíodos de signo diferente?}]

\subsection{Estimación del modelo CAPM}
\label{subsec:res_capm}

[\textit{Presentar resultados de la regresión CAPM para las carteras o activos analizados.}]

\subsection{Estimación del modelo Fama-French de tres factores}
\label{subsec:res_ff3f}

[\textit{Presentar resultados del modelo FF3F.}]

El Cuadro~\ref{tab:ff3f} recoge las estimaciones MCO del modelo FF3F para las
carteras analizadas.

\begin{table}[H]
  \centering
  \caption{Estimación del modelo Fama-French de tres factores}
  \label{tab:ff3f}
  \begin{tabular}{lcccccc}
    \toprule
    Cartera & $\hat\alpha$ & $\hat\beta$ & $\hat{s}$ & $\hat{h}$ & $\bar{R}^2$ & $N$ \\
    \midrule
    [\textit{Cartera 1}] & & & & & & \\
    [\textit{Cartera 2}] & & & & & & \\
    [\textit{...}]       & & & & & & \\
    \bottomrule
  \end{tabular}
  \begin{minipage}{\linewidth}
    \small\vspace{0.3em}
    \textit{Nota:} Errores estándar Newey-West entre paréntesis.
    $^{*}p<0.10$, $^{**}p<0.05$, $^{***}p<0.01$.
  \end{minipage}
\end{table}

\subsection{Estimación del modelo de cuatro factores con el factor ESG}
\label{subsec:res_ff4f}

[\textit{Este es el resultado central del trabajo. Presentar y discutir.}]

El Cuadro~\ref{tab:ff4f} recoge los resultados de la estimación del modelo
de cuatro factores (FF4F-ESG).

\begin{table}[H]
  \centering
  \caption{Estimación del modelo de cuatro factores FF4F-ESG}
  \label{tab:ff4f}
  \begin{tabular}{lccccccc}
    \toprule
    Cartera & $\hat\alpha$ & $\hat\beta$ & $\hat{s}$ & $\hat{h}$ & $\hat{e}$ (HEL) & $\bar{R}^2$ & $N$ \\
    \midrule
    [\textit{Cartera 1}] & & & & & & & \\
    [\textit{Cartera 2}] & & & & & & & \\
    [\textit{...}]       & & & & & & & \\
    \midrule
    \textbf{HEL (factor)} & & & & & & & \\
    \bottomrule
  \end{tabular}
  \begin{minipage}{\linewidth}
    \small\vspace{0.3em}
    \textit{Nota:} Errores estándar Newey-West entre paréntesis.
    $^{*}p<0.10$, $^{**}p<0.05$, $^{***}p<0.01$.
    HEL = \textit{High ESG minus Low ESG}.
  \end{minipage}
\end{table}

[\textit{Interpretar: ¿Es $\hat{e}$ significativo? ¿Cuál es su signo?
¿Mejora el $\bar{R}^2$ con respecto al FF3F? ¿Se reducen los alfas?}]

\subsection{Comparación de modelos y test GRS}
\label{subsec:comparacion}

[\textit{Comparar los tres modelos (CAPM, FF3F, FF4F-ESG) usando $\bar{R}^2$,
criterios AIC/BIC y el test GRS.}]

El Cuadro~\ref{tab:comparacion} resume los criterios de comparación entre modelos.

\begin{table}[H]
  \centering
  \caption{Comparación de modelos}
  \label{tab:comparacion}
  \begin{tabular}{lcccc}
    \toprule
    & CAPM & FF3F & FF4F-ESG \\
    \midrule
    $\bar{R}^2$ medio   & & & \\
    $|\hat\alpha|$ medio & & & \\
    GRS $F$-stat        & & & \\
    GRS $p$-value       & & & \\
    AIC medio           & & & \\
    \bottomrule
  \end{tabular}
\end{table}

\subsection{Análisis de robustez}
\label{subsec:robustez}

[\textit{Describir los tests de robustez realizados:
¿subperíodos? ¿diferentes puntos de corte ESG? ¿diferentes fuentes de datos?
¿ponderación igual vs. por capitalización?}]

% ═════════════════════════════════════════════════════════════════════════════
\section{Conclusiones}
\label{sec:conclusiones}
% ═════════════════════════════════════════════════════════════════════════════

[\textit{Resumir hallazgos principales, implicaciones y limitaciones.}]

Este trabajo ha investigado si los criterios ESG generan una prima de riesgo
estadísticamente significativa en el mercado [\textit{especificar}], incorporando
un factor HEL al modelo de tres factores de Fama y French (1993).

Los principales hallazgos son los siguientes:

\begin{enumerate}
  \item [\textit{Hallazgo 1: significatividad del factor HEL.}]
  \item [\textit{Hallazgo 2: dirección del efecto (prima positiva/negativa/nula).}]
  \item [\textit{Hallazgo 3: mejora del ajuste del modelo.}]
  \item [\textit{Hallazgo 4: robustez.}]
\end{enumerate}

Estos resultados son [\textit{consistentes / inconsistentes}] con la hipótesis de
Pástor, Stambaugh y Taylor (2021, 2022) de que los inversores con preferencias
ESG aceptan menores rendimientos financieros, y con la hipótesis de exclusión
de Hong y Kacperczyk (2009).

Las principales limitaciones del trabajo son: [\textit{tamaño de muestra,
disponibilidad de datos ESG, período relativamente corto, posible sesgo de
supervivencia, heterogeneidad en puntuaciones ESG entre proveedores.}]

Como líneas de investigación futuras se propone: [\textit{ampliar el período,
comparar entre mercados, usar distintas fuentes ESG, incorporar el factor
momentum.}]

% ─── Bibliografía ────────────────────────────────────────────────────────────
\newpage
\printbibliography[title={Referencias Bibliográficas}]

% ─── Apéndices ───────────────────────────────────────────────────────────────
\appendix

\section{Detalle de la construcción de carteras}
\label{app:carteras}

[\textit{Describir con detalle el proceso de construcción de las carteras.}]

\section{Resultados adicionales de robustez}
\label{app:robustez}

[\textit{Tablas y figuras adicionales.}]

\end{document}
